\scanpage{53}
\mspage{- 13 -}

n$^{\circ}$385

7. \uline{Produit cartésien de deux distributions}.

Soient $X$ et $Y$ deux variétés de classe $C^r$, $a$ un point de $X$ et $b$ un point de $Y$.

Il existe une application linéaire $\iota$ de $D_a^r(X)\otimes D_b^r(Y)$ dans l'espace $B$ des formes bilinéaires sur $C_a^r(X)\times C_b^r(Y)$ caractérisée par la formule :
\[
\iota(\mu\otimes\unsure{\nu|\mathfrak{d}})(f,g)=\langle\mu,f\rangle\,\langle\nu,g\rangle.\tag*{(31)}
\]

Par ailleurs, on définit une application linéaire $\iota'$ de $D_{a,b}^r(X\times Y)$ dans $B$ par la formule :
\[
\iota'(\pi)(f,g)=\langle\pi,f\times g\rangle.\tag*{(32)}
\]

PROPOSITION 4.- \uline{Il existe une application linéaire} $\rho_{a,b}$ \uline{de} $D_{a,b}^r(X\times Y)$ \uline{dans} $D_a^r(X)\otimes D_b^r(Y)$ \uline{et une seule telle que} $\iota'=\iota\circ\rho_{a,b}$. \uline{L'application} $\rho_{a,b}$ \uline{est injective et son image est le sous-espace de} $D_a^r(X)\otimes D_b^r(Y)$ \uline{engendré par les éléments} $\mu\otimes\nu$ \uline{avec} $v(\mu)+v(\nu)\leq r$. \uline{En particulier}, \uline{si} $r$ \uline{n'est pas un entier}, \uline{l'application} $\rho_{a,b}$ \uline{est bijective}.

Nous introduirons un système de coordonnées $x_1,\ldots,x_n$ sur $X$ au voisinage de $a$ et tel que $x_i(a)=0$ pour tout $i$ ; de même, on introduit un système de coordonnées $y_1,\ldots,y_m$ sur $Y$ au voisinage de $b$ tel que $y_j(b)=0$ pour tout $j$. D'après la prop.2 il existe des bases $\{\Delta_a^\alpha\}$ de $D_a^r(X)$ et $\{\Delta_b^\beta\}$ de $D_b^r(Y)$ caractérisées par les formules :
\[
\langle\Delta_a^\alpha,x^{\alpha'}\rangle=\delta_{\alpha\alpha'}\qquad
\langle\Delta_b^\beta,y^{\beta'}\rangle=\delta_{\beta\beta'}\tag*{(33)}
\]
pour $\alpha$ et $\alpha'$ dans $N_r^n$ et $\beta$ et $\beta'$ dans $N_r^m$. Par conséquent, le produit tensoriel $D_a^r(X)\otimes D_b^r(Y)$ admet pour base les éléments $\Delta_a^\alpha\otimes\Delta_b^\beta$ pour $\alpha$ dans $N_r^n$ et $\beta$ dans $N_r^m$.

Par ailleurs (référence), on définit un système de coordonnées $z_1,\ldots,z_{n+m}$ sur $X\times Y$ au voisinage de $(a,b)$ par les conditions :
% archive sheet number: 52
